\documentclass[12pt]{article}
\usepackage{amsmath}
\usepackage{graphicx}
\usepackage{hyperref}

\title{Social Justice as a Neural Network System}
\author{IYR \and ChatGPT}
\date{}

\begin{document}

\maketitle

\section*{Problem Statement}

This work models social justice as a neural network system, where societal actors, policies, and events interact as nodes and edges. The goal is to understand how fairness, bias, and systemic outcomes emerge from distributed interactions.

\section*{Repository Reference}

The implementation corresponding to this chapter is available at:

\begin{quote}
\url{https://github.com/spacetracker-collab/SocialJusticeNeuralNetwork.git}
\end{quote}

This repository provides a neural modeling framework where social variables are encoded and processed through layered architectures to simulate systemic outcomes.

\section*{Relation to Existing Work}

The proposed system aligns with prior research in social and graph-based neural modeling:

\begin{itemize}
\item Social interaction models such as SocialNet demonstrate how pairwise and aggregate interactions can predict collective behavior \cite{heras2019socialnet}.
\item Graph-based neural approaches like Social-STGCNN model interactions as structured networks and show improved predictive performance in social trajectory tasks.
\item Large-scale social network datasets, such as GitHub social graphs, represent individuals as nodes and relationships as edges, supporting neural interpretations of social systems.
\end{itemize}

These works collectively support the thesis that:
\begin{quote}
Social systems can be effectively modeled as neural or graph-based computational structures.
\end{quote}

\section*{Model Description}

\subsection*{Architecture}

The model is structured as a multi-layer neural network:

\begin{itemize}
\item Input Layer: Encodes social variables (actors, policies, states)
\item Hidden Layers: Capture interactions and dependencies
\item Output Layer: Represents outcomes (equity distribution, bias metrics)
\end{itemize}

\subsection*{Mathematical Representation}

\[
y = f(Wx + b)
\]

Where:
\begin{itemize}
\item $x$ = encoded social state
\item $W$ = interaction weights
\item $b$ = bias terms
\item $f$ = nonlinear activation
\end{itemize}

\section*{Neural Mapping}

\begin{itemize}
\item Individuals / Institutions $\rightarrow$ Neurons
\item Social Interactions $\rightarrow$ Weighted edges
\item Policies $\rightarrow$ Parameter shifts
\item Outcomes $\rightarrow$ Network activations
\end{itemize}

\section*{Training Process}

The model is trained on structured or synthetic datasets representing:

\begin{itemize}
\item Social interactions
\item Policy interventions
\item System responses
\end{itemize}

Training involves iterative optimization and loss minimization.

\section*{Results}

Example output:

\[
[0.2967,\ 0.3981,\ 0.3051]
\]

This represents probabilistic system states corresponding to different social outcomes.

\section*{Inference and Insights}

Key insights include:

\begin{itemize}
\item Social systems exhibit distributed intelligence rather than centralized control.
\item Bias emerges from structural weights rather than isolated actors.
\item Small perturbations can propagate nonlinearly across the system.
\item Emergent behavior dominates deterministic prediction.
\end{itemize}

\section*{Case Studies}

\subsection*{Policy Intervention}

Adjusting weights corresponding to policy nodes alters system equilibrium.

\subsection*{Network Influence}

Highly connected nodes disproportionately affect global outcomes.

\section*{Sports Theory Mapping}

\begin{itemize}
\item Teams $\rightarrow$ Neural clusters
\item Players $\rightarrow$ Nodes
\item Strategy $\rightarrow$ Weight optimization
\item Match outcomes $\rightarrow$ Activation outputs
\end{itemize}

This reinforces the idea that competitive systems behave like neural optimization processes.

\section*{Extensions}

Future work includes:

\begin{itemize}
\item Multi-agent reinforcement learning
\item Real-time data integration
\item Global-scale system modeling
\item Meta-learning (AI modeling AI systems)
\end{itemize}

\section*{Conclusion}

This work demonstrates that social justice systems can be interpreted as neural networks, enabling a unified computational framework for analyzing fairness, bias, and systemic dynamics.

\section*{References}

\begin{thebibliography}{9}

\bibitem{repo}
SocialJusticeNeuralNetwork Repository.  
\url{https://github.com/spacetracker-collab/SocialJusticeNeuralNetwork.git}

\bibitem{heras2019socialnet}
Heras, F. J. H., Romero-Ferrero, F., Hinz, R. C., de Polavieja, G. G.  
``Deep attention networks reveal the rules of collective motion in zebrafish.''  
PLOS Computational Biology, 2019. :contentReference[oaicite:0]{index=0}

\bibitem{socialgraph}
Rozemberczki, B., Allen, C., Sarkar, R.  
``Multi-scale Attributed Node Embedding.''  
GitHub Social Network Dataset, 2019. :contentReference[oaicite:1]{index=1}

\bibitem{stgcnn}
Mohamed, A., Qian, K., Elhoseiny, M., Claudel, C.  
``Social-STGCNN: Social Spatio-Temporal Graph Convolutional Neural Network.''  
arXiv, 2020. :contentReference[oaicite:2]{index=2}

\end{thebibliography}

\end{document}
